\section{Evaluation \& Results}
\label{sec:eval}

\subsection{Setup, metrics, and staging}
\label{sec:eval-setup}

All numbers are measured on the held-out benchmark of the studied project (Tier-3; $n=60$ cases / 59
elements, of which 35 are addressable fillers). The split is region-disjoint by construction, but a
data audit found \textbf{element-level leakage}: 12 of the 59 held-out target elements (9 fillers, 3
walls) also appear as training targets under a different render, because an element can sit in two
regions. We disclose this directly and note that it cannot affect the headline results by
construction: the oracle diagnostics (\secref{sec:eval-oracle}, \secref{sec:eval-depth}) read only
the IFC graph and learn nothing, and the realized position-slot result (\secref{sec:eval-realized})
uses a \emph{deterministic} OpenCV detector with no trained weights --- both are leakage-proof. The
only learned component touched by the split is the per-field VLM profile
(\secref{sec:eval-neural}); re-profiling on the leakage-clean 48-case subset leaves it unchanged
(storey/class $100\%$, position-slot/size $0\%$), so the architectural conclusion does not depend on
the leaked cases. We structure the
evaluation in stages: first, an \emph{oracle} experiment establishing the symbolic-layer ceiling under
perfect extraction (\secref{sec:eval-oracle}); second, neural extraction accuracy per field
(\secref{sec:eval-neural}); third, how noisy signals combine without destroying recall
(\secref{sec:eval-union}); fourth, the realized extractor with calibrated, selectively-predictive
routing (\secref{sec:eval-realized}); fifth, the depth analysis (\secref{sec:eval-depth}); and last,
failure decomposition and workflow-facing cost (\secref{sec:eval-failure}). This staging separates
symbolic retrieval quality from neural extraction quality --- the two failure modes require different
fixes, and conflating them produces misleading aggregate metrics. We report \textbf{GT-in-Pool} as the
primary retrieval metric rather than Top-1, because Top-1 conflates retrieval failure (element never
in the pool) with ranking failure (element in the pool but ranked poorly); guaranteed retrievability is
the architectural contribution, and ranking is where the spatial address acts. Table~\ref{tab:main} consolidates the headline numbers; the remainder of
this section derives and explains each block.

\begin{table}[ht]
  \centering
  \small
  \begin{tabular}{lrrrr}
    \toprule
    Method & GT-in-Pool & Top-1 & Top-10 & MRR@10 \\
    \midrule
    \multicolumn{5}{l}{\emph{External baselines, full held-out pool ($n=60$)}} \\
    \quad Lexical (BM25)                          & 100 & 1.7  & 15.6 & 0.058 \\
    \quad Dense (MiniLM)                          & 100 & 1.7  & 25.0 & 0.108 \\
    \quad Zero-shot VLM --- full pool             & 100 & 3.3  & 15.0 & 0.086 \\
    \quad Zero-shot VLM --- sibling shortlist     & 100 & 3.3  & 20.0 & 0.097 \\
    \midrule
    \multicolumn{5}{l}{\emph{Our pipeline, realized end-to-end ($n=60$)}} \\
    \quad G8 extraction $\to$ deterministic retrieval & ${\sim}100$ & 6.7 & 30.0 & 0.110 \\
    \midrule
    \multicolumn{5}{l}{\emph{Oracle ceiling, same pools ($n=60$)}} \\
    \quad Coarse: storey $+$ class                & 100 & 4.9  & 31.5 & 0.137 \\
    \quad \quad $+$ \texttt{object\_type}         & 100 & ---  & 76.0 & ---   \\
    \quad \quad $+$ structured spatial address    & 100 & \textbf{78.5} & \textbf{98.1} & \textbf{0.854} \\
    \midrule
    \multicolumn{5}{l}{\emph{Realized address, addressable fillers ($n=35$)}} \\
    \quad Modal-prior slot (floor)                & 100 & 6.6  & ---  & ---   \\
    \quad OpenCV slot $+$ soft rerank             & 100 & \textbf{67.6}$_{[53.6,81.6]}$ & 80.9 & --- \\
    \bottomrule
  \end{tabular}
  \caption{Consolidated held-out results. Every method ranks the \emph{same} retrieved pools with the
  ground truth in-pool throughout (GT-in-Pool, \%), so the table isolates \emph{ranking}, not
  retrieval. Top-1/Top-10/MRR in \%. External text retrieval and a zero-shot Qwen2.5-VL-7B reranker
  (base model, no adapter) stay at single-digit Top-1 --- the VLM at chance --- while the oracle
  structured address reorders the identical pool to 78.5\%. The realized filler block shows how much
  survives a real detector: one deterministic specialist lifts Top-1 from the 6.6\% floor to 67.6\%
  (bootstrap 95\% CI in subscript); deferring the least-confident ${\sim}20\%$ further lifts the
  answered subset to 80.6\% (\secref{sec:eval-realized}). Dashes mark quantities not measured for that
  row.}
  \label{tab:main}
\end{table}

\subsection{Oracle experiment: the symbolic ceiling, and where discrimination lives}
\label{sec:eval-oracle}

Under perfect extraction the symbolic layer retains the target in \textbf{100\% of cases}
(GT-in-Pool) and compresses the live pool along the fingerprint ladder --- from $1{,}233$
elements (whole-model) to 46 (storey+class) toward single digits as spatial constraints are added. The architecture is therefore provably sound; the question is where,
inside the pool, the discriminating information lives.

\paragraph{The coarse floor saturates.} Restricting to the target's storey and IFC class leaves a
median confusable set of 46 (mean 112) and an oracle Top-1 of only \textbf{4.9\%} --- \emph{below} the
realized end-to-end reranker (6.7\%) and equal to it on Top-10 (31.5\% vs.\ 30.0\%). Ontology is
necessary but cannot separate same-class siblings: it is the floor, not the discriminator. Adding the
richest attribute (the family/type string \texttt{object\_type}) shrinks the confusable set to a
median of 13 ($3.8\times$) and lifts oracle Top-10 to 76\%, but uniquely identifies only 2 of 60
targets. What remains is the irreducibly \emph{relational} residue.

\paragraph{The type-conditional address closes it.} Supplied at an oracle level, the class-conditional
address of \secref{sec:address} takes the real retrieved pool from a coarse Top-1 of 4.9\% to
\textbf{78.5\%}, and Top-10 from 31.5\% to \textbf{98.1\%} (MRR $0.137 \to 0.854$), at zero recall
cost. The gain decomposes exactly along the type-conditional split --- fillers reach Top-1
\textbf{91.0\%} (position-slot), walls \textbf{64.2\%} (connectivity fingerprint; within same-storey
walls the fingerprint collapses a median confusable set of $110 \to 2$, uniquely identifying 10 of 22
walls where \texttt{object\_type} identifies none). No uniform descriptor achieves this: the
position-slot is meaningless for a wall, the fingerprint degenerate for a window. Generic retrieval
baselines confirm the gap is not an artifact of our own pipeline. Lexical and dense retrievers plateau
at Top-1 1.7\% (Top-10 15.6\% / 25.0\%) on the same pools, and a \emph{zero-shot} Qwen2.5-VL-7B
reranker --- the same backbone our extractor fine-tunes, given the marked site photo, the marked
floorplan, and the candidate pool --- stays at chance: Top-1 3.3\% over the full pool and 2.9\% on
fillers even when the pool is pre-filtered to the same-storey, same-class confusable set
(per-case chance 1.8--2.4\%; candidates are shuffled so that copying the list order scores at chance,
which is what the model does). On those same pools the oracle spatial address reaches 78.5\%
(Figure~\ref{fig:external}). A strong general VLM cannot recover the relational slot from the image
alone; the structured address, not the backbone, is what discriminates identical siblings.

\begin{figure}[ht]
  \centering
  \includegraphics[width=0.85\textwidth]{external_baseline.png}
  \caption{Baselines against the oracle ceiling on identical pools. Lexical and dense text retrieval,
  a zero-shot Qwen2.5-VL-7B reranker (base model, no adapter; both the full pool and the
  same-storey/class sibling shortlist), and the deployed G8 pipeline all plateau in single-digit
  Top-1 --- the zero-shot VLM at chance --- with the ground truth in-pool throughout; the oracle
  spatial address reorders the same pool to 78.5\%/98.1\%. The bottleneck is ranking inside the pool,
  and the missing signal is the relational address.}
  \label{fig:external}
\end{figure}

\subsection{Neural extraction quality: what the fine-tuned VLM does and does not recover}
\label{sec:eval-neural}

Per-field evaluation shows fine-tuning works where it was aimed: against zero-shot prompting
(\texttt{ifc\_class} 63.3\%, storey 30.0\%), the fine-tuned models saturate both at \textbf{100\%},
reach predicate slot-accuracy 82.8--86.2\% versus 43.1\%, and lift relation direction to 82.1\% versus
${\sim}0\%$. End-to-end, however, the best configuration delivers
Top-10 30.0\% / Top-1 6.7\% / MRR@10 0.110 --- with the target in the pool ${\sim}100\%$ of the time.
The decisive diagnostic is the held-out per-field re-profile (Figure~\ref{fig:vlm}): the VLM is
saturated on the \emph{coarse prefix} (storey, class --- exactly the fields the oracle shows are
non-discriminating) and recovers the \emph{discriminating} structured fields at \textbf{0\%} (the
position-slot and the size band), and emits a relation-direction field in 57\% of held-out cases
(this 57\% is an \emph{emission rate} --- whether any direction is produced --- not the 82\%
direction \emph{accuracy} reported above on the thesis evaluation set; the two numbers measure
different quantities and are not comparable). The conclusion is
architectural, not ``fine-tune harder'': delegate the slot and size to deterministic visual
specialists, keep the VLM on the coarse prefix and relation typing where it is reliable. This is the
neuro-symbolic interface in practice --- learn where the network is reliable, compute
deterministically where it is not \citep{mao2019nscl}.

\begin{figure}[ht]
  \centering
  \includegraphics[width=0.85\textwidth]{vlm_profile.png}
  \caption{Perception takeaway: per-field re-profile of the fine-tuned VLM on the held-out set. The
  coarse prefix (storey, class) saturates; the discriminating fields (position-slot, size) are
  recovered at 0\% --- motivating delegation to deterministic specialists.}
  \label{fig:vlm}
\end{figure}

\subsection{Combining signals: union, never intersection}
\label{sec:eval-union}

How should a mature signal (storey+type) combine with an immature one (spatial topology)? The strategy
ablation is unambiguous: the \textbf{union} of constraint sets is the only strategy that preserves
full GT-in-Pool while still gaining the spatial constraints' compression and ranking benefit; hard
intersection over-prunes under noisy extraction, and spatial-only is too brittle. The measurement says why: treating each extracted field as a hard
filter multiplies per-field recalls --- with measured single-field reliabilities (storey
$\approx 0.66$, class $\approx 0.50$), a hard conjunction over the attribute fields collapses joint
recall toward $\prod r \approx 0.009$. The filter evicts the ground truth far more often than it
removes a distractor. The resolution is to relocate determinism: \textbf{the executor is deterministic
given the structured record; the ranking is a calibrated soft prior inside a recall-fixed pool.} The
address never removes a candidate; it re-weights one. Recall stays at 100\% by construction, and the
unreliable signal is spent only on ordering. As a design principle this generalizes: when combining
mature and immature signals, union --- intersection amplifies extraction error and destroys recall.

\subsection{Realizing the address: one specialist, calibrated, selectively predictive}
\label{sec:eval-realized}

The oracle ceiling is a statement about information; this section measures how much survives a real,
imperfect detector. The honest floor is stark: the fine-tuned reranker emits position as free text and
recovers a usable slot in \textbf{0 of 35} held-out fillers; even a modal-prior slot guess reaches
only 6.6\%. We report four results.

\paragraph{(1) One realizable extractor closes most of the gap.} The OpenCV position-slot detector
recovers the opening count exactly in 83\% of cases (29/35) and the full slot in 74\% (joint
$(i,M)$). Fed into the soft rerank, it lifts filler Top-1 from the 6.6\% floor to \textbf{67.6\%}
(bootstrap 95\% CI $[53.6, 81.6]$, $n=35$) and Top-10 to 80.9\%, at zero recall cost. The residual
error concentrates in the ordering index, not in the perception of the openings. The wide interval is
a direct consequence of the small addressable set and is reported, not smoothed over. The detector's
three pixel thresholds are not a knife-edge fit: sweeping each $\pm25\%$ leaves realized Top-1 in
$[55.1, 83.1]$ (never near the 6.6\% floor), two of the three are inert, and only the same-wall
perpendicular band matters --- and the deployed value is mid-range, not tuned to the peak.

\paragraph{(2) A representational prerequisite: the address must be image-recoverable.} The
ground-truth slot in the source data is numbered along each wall's IFC local-X axis --- a modelling
artefact invisible in any image. Scored against that convention, detector and ground truth disagree on
16 of 35 fillers, \emph{all} exact mirrors $i \mapsto M{-}1{-}i$; under an image-coordinate convention
the disagreement vanishes and joint accuracy rises from an apparent 34\% to the true 74\%. An address
field is realizable only if it is defined in a frame the evidence carries --- a substantive design
constraint, not bookkeeping. The same constraint explains a negative result: the wall fingerprint's
load-bearing fields depend on IFC wall-\emph{instance} boundaries that collinear instances erase in
the floorplan poch\'e (length-band recovery 5/17, realized Top-1 $\approx$ floor), so the wall address
is oracle-discriminative but \emph{not} image-realizable, and realization is scoped to the filler
slot.

\paragraph{(3) The confidence is calibratable --- and reweighting is a measured no-op.} Routing is
only legitimate if confidence tracks correctness \citep{guo2017calibration}. On the held-out fillers
the raw detector confidence is positively discriminative (AUROC \textbf{0.80}, bootstrap 95\% CI
$[0.58, 0.99]$) and moderately mis-calibrated (ECE 0.206, 95\% CI $[0.09, 0.32]$); temperature
scaling reduces it to 0.172 (Figure~\ref{fig:caldiag}). The intervals are wide --- at $n=35$ the
calibration evidence is suggestive, not conclusive, and we treat the gate as a measured operating
point rather than a precise estimate. But
the calibrated weight does \emph{not} pay off as a rerank weight: because the slot is the finest
tiebreaker inside a storey$\times$class bucket, any strictly-positive weight induces the identical
ordering --- hard match, raw weight, and calibrated weight all yield the same 67.6\%. We report this
no-op rather than bury it.

\paragraph{(4) The payoff is selective prediction.} The calibrated confidence pays off as a deferral
threshold: sweeping it traces a coverage--accuracy curve (Figure~\ref{fig:rerank}), and deferring the
least-confident ${\sim}20\%$ of cases lifts Top-1 on the answered subset from 67.6\% to
\textbf{80.6\%} \citep{geifman2017selective}. A worked case (Figure~\ref{fig:demo}) shows the
mechanism: the detector predicts slot 1 of 10 where the truth is 8 of 10 --- an error --- but its
calibrated confidence (0.05) is below threshold, so the case routes to \emph{defer} and surfaces the
candidates instead of a confident wrong GUID. For a triage tool with downstream consequence, this is
the correct behavior, and the coverage--accuracy curve, not a point estimate, is the honest report.

\begin{figure}[ht]
  \centering
  \includegraphics[width=0.85\textwidth]{calibration_diag.png}
  \caption{Calibration gate. Raw detector confidence is positively discriminative (AUROC 0.80) and
  moderately mis-calibrated (ECE 0.206); temperature scaling reduces ECE to 0.172.}
  \label{fig:caldiag}
\end{figure}

\begin{figure}[ht]
  \centering
  \includegraphics[width=0.85\textwidth]{calibrate_rerank.png}
  \caption{Selective prediction. Sweeping the deferral threshold traces a coverage--accuracy curve;
  deferring the least-confident ${\sim}20\%$ lifts answered-set Top-1 from 67.6\% to 80.6\%.}
  \label{fig:rerank}
\end{figure}

\begin{figure}[ht]
  \centering
  \includegraphics[width=\textwidth]{case_AP_SK_092.png}
  \caption{A worked DEFER case. The detector predicts slot 1 of 10 where the truth is 8 of 10, but
  its calibrated confidence (0.05) is below threshold: the case routes to defer and surfaces the
  candidate pool rather than emitting a confident wrong GUID.}
  \label{fig:demo}
\end{figure}

\subsection{The depth law: relational context saturates at one hop}
\label{sec:eval-depth}

Where should relational context live --- extracted deep at inference, or compiled into the node? Under
an oracle, deeper context keeps shrinking the confusable set (a Weisfeiler--Leman-style refinement
toward singletons). But the relations must be recovered from an image, and per-hop extraction
reliability falls as $0.40 \to 0.05 \to 0$, so the \emph{realizable} confusable set saturates
immediately (Table~\ref{tab:depth}, Figure~\ref{fig:depth}): median 13 (attributes) $\to$
\textbf{8.2} at one hop $\to$ 8.1 at two and three. This is the depth law --- deeper context is
informative but not recoverable --- and its failure mode is the documented error compounding of
multi-hop reasoning over knowledge graphs \citep{chakraborty2024multihop}. The training side
corroborates it: supervising depth-3 chains yielded ${\approx}0$ realizable improvement and \emph{cost}
capacity --- under a fixed LoRA rank, depth-$\ge 2$ chains degraded the coarse \texttt{ifc\_class}
recovery, and the deep-chain extraction itself ran at ${\sim}5\%$ and hurt end-to-end performance by
${\sim}13$ points. The prescription is subtractive: stop training deep
chains, compile depth into the node, extract at depth $\le 1$.

\begin{table}[ht]
  \centering
  \begin{tabular}{lc}
    \toprule
    relational depth & realizable median $|C|$ \\
    \midrule
    attributes only (depth 0) & 13 \\
    \textbf{+ 1 hop} & \textbf{8.2} \\
    + 2 hops & 8.1 \\
    + 3 hops & 8.1 \\
    \bottomrule
  \end{tabular}
  \caption{Realizable median confusable-set size by relational depth (held-out set). All realizable
  discrimination is captured at one hop.}
  \label{tab:depth}
\end{table}

\begin{figure}[ht]
  \centering
  \includegraphics[width=0.85\textwidth]{depth_saturation.png}
  \caption{Information versus realizability. The oracle confusable set keeps shrinking with depth; the
  reliability-weighted realizable set saturates at one hop --- the depth law.}
  \label{fig:depth}
\end{figure}

\subsection{Failure decomposition, latency, and triage effort}
\label{sec:eval-failure}

Because the pipeline localizes errors to named fields rather than opaque scores, the residual
end-to-end loss decomposes cleanly: it is dominated by \emph{extraction} errors --- missing or
incorrect spatial predicates and the lowest-accuracy fields (subtype, material, position-context) ---
not by retrieval logic. Every failure has a traceable stage, which is the operational
differentiator over black-box matching: the roadmap is ``improve the named weak fields,'' not
``redesign the planner.'' Latency favors the constrained pipeline (${\sim}1$\,s vs.\ ${\sim}4.5$\,s
for the ReAct agent, with exact repeatability). The workflow-facing cost metric makes the ceiling
concrete (Figure~\ref{fig:triage}): the expected number of candidates a coordinator must inspect falls
from a median 38 under manual scan to 0.5 under a perfect address --- with 15 seconds per inspection an
explicit proxy assumption, not a user-study result.

\begin{figure}[ht]
  \centering
  \includegraphics[width=0.8\textwidth]{triage_effort.png}
  \caption{Triage-effort proxy on the same held-out pools. Expected inspections per case fall from a
  median 38 (manual scan) to 0.5 (perfect address); the realized specialist sits between, with
  success@1 78.5\% under the oracle address.}
  \label{fig:triage}
\end{figure}
